\documentclass{article}
\usepackage{amsmath}
\newcommand{\argmin}{\arg\!\min} 
\newcommand{\argmax}{\arg\!\max} 
\usepackage{amssymb}
\usepackage{amsfonts}
\usepackage[T1]{fontenc}
\usepackage{bm}
\usepackage{array}
\usepackage{graphicx}
\usepackage[utf8]{inputenc}

\usepackage[   
	backend=biber,
	style=alphabetic,
	sorting=ynt 
]{biblatex}
\addbibresource{references.bib}

\title{Geometric redundancy resolution in inverse kinematics}
\date{\today}
\author{Marko Guberina}

\begin{document}
\maketitle


\section{Introduction}
The concept of manipulability has been around since 1985.
The idea is very straightforward. Take a ball of joint velocities in joint space,
and map it, via the Jacobian, to operational space. The output is an ellipsoid.
This ellipsoid tells you how easy it is to move in a particular direction,
or, equivalently, how sensitive you are to forces in a particular direction.
For the most part, it was used only to judge the quality of a particular posture.
If you have a 6dof robot, there is no need to do anything else,
as you can simply select the IK solution with the highest score.
But, if you have a redundant robot, there are whole dofs to explore.
However, in the vast majority of cases, 
the focus was on the volume of the ellipsoid,
or on the length of the shortest axis (as at singularity the shortest axis of the
ellipsoid is 0 in length) [TODO: have some citations about 
the pre-geometric use of MEs in redundant robots].
There were some minor efforts in designing a target manipulability ellipsoid (ME)
\cite{lee1989dual,lee2016humanoid},
but nothing came of that as far as I'm aware,
as computing and setting up task compatibility is extremely time-consuming
and tedious.

A turning point came with \cite{jaquier2021geometry} when the authors
finally introduce differential geometry to manipulability ellipsoids.
We've known about screw theory and the benefits of interpreting
rigid body transformations as Lie group elements since at least the 90s,
but noone bothered to extend the geometric approach to manipulability before.
With this we can actual track prescribed manipulability ellipsoids,
instead of just computing them.

So far there have been only 2 (and a half) applications of 
manipulability tracking.
The first is in learning from demostrations, specifically from a human.
Humans can do most tasks we care about, but they have a different
kinematic structure than the target robot.
While you can track the end-effector position, you can't do anything
with just joint angles (as they don't map to the robot), 
and you lose all information about
posture, which is critical when you need to exert high forces for example.
But you can find the human's manipulability ellipsoid, 
and then try to achieve that
as it is independent from the underlying kinematic structure!
Well, not exactly, but there is a mapping between the human's ME 
and the robot's ME \cite{reithmeir2022human}.
The second one is singularity avoidance, for example \cite{maric2021riemannian}.
The third one isn't even manipulability tracking, just manipulability
index (ellipsoid volume) maximization.


Ideally, we don't need demonstrations. Instead, we build an algorithm
which can generate a solution for the given motion problem from first 
principles, and then execute the solution --- 
all you had to do was specify the task.
For example, there are many tasks in which we know which path we want the robot 
to follow,
ex. welding, polishing etc.

Moreover, sometimes we just want to move the robot a bit.
As it turns out redundancy resolution is apsolutely critical 
when there are left-over dofs to exploit. In fact, 
the higher the nullspace dimension, the more important 
it is to do some kind of reduncancy resolution.

Let's look at this case a bit more carefully: there is a good
prescribed path, and we want to track it.
To keep things simple for now, let's say there are no 
external contacts and forces, i.e. the path is in free-space.
The existing redundancy resolution methods don't
really offer anything to directly benefit in this task.
As noted in \cite{lee1989dual} [TODO: or wherever idk now],
manipulability maximization can actually hurt performance.
Singularity avoidance helps you avoid singularities,
but that's it. Conceivably, one should be able to design
a useful and feasible redundancy resolution method,
which would then avoid singularities as a byproduct.
Thus, it is sensible to try to design manipulability ellipsoids
strictly from the path geometry, and use them for 
redundancy resolution.

This was done in \cite{yan2021decentralized}.


\subsection{Why is redundacy resolution still useful?}
\paragraph{Isn't this magically solved during optimization?}
Here one cound say, ah, but we have trajectory optimization,
so what is the point of crafting these additional structures 
when you could just throw in all your costs and constraints
into an optimal control formulation?
Depending on how the OCP is constructed, that might be the case.
However:
\begin{itemize}
\item in general these OCPs are nonconvex, so you're not actually
	finding the optimal solution
\item MPC with manipulability maximization has higher manipulability
	than MPC without it (TODO: prove by demonstration),
	so you can (and should) put it as a conditioning cost
\item you can put this in your traj. opt. and get a better
	conditioned problem, allowing you to add more costs/constraints
\item even if traj. opt. did somehow leverage manipulability,
	it would be good to actually show this - more understanding,
	more profit
\end{itemize}





The goal of is to leverage the manipulability tracking
controllers described in \cite{jaquier2018geometry}.

All the relevant mathematics can be found in the paper text, although a reader 
without experience with differential geometry
will most likely not be able to parse the finer details.
Hence this text will only concern itself with methodology.

\section{Manipulability tracking as a QP cost}
Throughout the text, we restrict ourselves to 3D linear velocities.
The methodology is identical for the full 6D case including angular velocities,
but the targets and effect on manipulator posture are more difficult to visualize and assess.
Hence $ v \in \mathbb{R}^{3}, J \in \mathbb{R}^{3 \times n}  $ and so on.
To cover both cases, we'll use $ d \in \left\{ 3,6 \right\}   $.
Assume there is a target manipulability ellipsoid $ \Lambda \in S^{ D }_{ ++ } \subset \mathbb{R}^{3 \times 3} $. 
The goal is to, for one instant in time, compute the vector of joint
velocities such that the current manipulability ellipsoid $ M  $.
This vector can then be used as the main kinematic task,
or a secondary task.
In the paper, the secondary task is achieved via direct nullspace projection.
However, this library (and myself) favour posing kinematic tasks
as quadratic programs (QPs) solved by appropriate QP solvers
as this framework makes it easier to combine different tasks,
and, more importantly, linear constraints.
Thus the very minor contribution of this work, apart from the implementation leveraging Pinocchio,
is the formulation of manipulability tracking as a quadratic cost.

\subsection{Derivation of target tangent vector in SPD space}
Utilizing the following inner product at $ T_{ M } S_{ ++ }^{ D }  $:
\begin{equation}
\left<\Sigma_{ 1 }, \Sigma_{ 2 } \right>_{ M }=
tr \left( 
M^{ -\frac{1}{2}  } \Sigma_{ 1 }
M^{ -1  } \Sigma_{ 2 }
M^{ -\frac{1}{2}  } 
\right) 
\end{equation}
as the Riemmanian norm, we can from the exponential and logarithmic maps
on the SPD manifold:
\begin{align}
\Lambda = \exp_{ M } L &=
M^{ \frac{1}{2}  } \exp (M^{ -\frac{1}{2}  }L M^{ -\frac{1}{2}  })M^{ \frac{1}{2}  }\\
L = \log_{ M } \Lambda &=
M^{ \frac{1}{2}  } \log (M^{ -\frac{1}{2}  }\Lambda M^{ -\frac{1}{2}  })M^{ \frac{1}{2}  }
\end{align}

Thus, provided the current manipulability ellipsoid $ M  $ and the target 
manipulability ellipsoid $ \Lambda  $, we get the tangent 
vector from $ M  $ to $ \Lambda  $
as $ L = \log_{ M } (\Lambda)  $.
Now we need to compute the vector of joint velocities $ v  $ such that
we follow $ L  $.

\subsection{Computing the derivative of the manipulability ellipsoid w.r.t. joint velocities}
We begin with the obvious computations:
\begin{gather}
		M = J J^{ T } \text{ definition of kinematic manipulability ellipsoid}\\
		\dot{M} = \dot{J} J^{ T } + J \dot{J}^{ T } \text{ applying } \frac{d }{d t}
\end{gather}
Here we set $ \dot{M} = L  $ to get the change in the manipulability ellipsoid we want.
To make this a QP cost, we need a quadratic of $ \dot{q}  $.
To get that, we need
$ \dot{M} = \mathcal{J} \times _{ 3 } \dot{q}^{ T }  $.
From this we can do:
\begin{gather}
		\min_{\dot{q}} \left\lVert 
\dot{M} - \mathcal{J} \times _{ 3 } \dot{q}^{ T }
		\right \rVert  \\
		=\min_{\dot{q}} \left\lVert 
vec(\dot{M}) - \mathcal{J}^{ T } _{ (3 )} \dot{q}
		\right \rVert  
\end{gather}
This is just rearranged linear algebra and it's better not to think about it too much
--- if the result is the same, it does not matter in which order one multiplies and sums.
But to make at least notation less confusing, $ vec (\cdot)  $ is matrix vectorization
(stack matrix rows into a single column vector (I think)),
and $ _{ (3) }  $ is tensor matricization along the 3rd index (same idea as vectorization,
the point is to write the tensor product as matrix vector multiplication and that's it).
The point here is we need to get this $ \mathcal{J}  $.
It can also be written like this:
\begin{equation}
\mathcal{J} = \frac{\partial J}{\partial q} \times _{ 2 } J +
\frac{\partial J}{\partial q} ^{ T } \times_{ 1 } J
\end{equation}

We need to look into $ \dot{J}  $ a bit, from the perspective of 
differentiating the forward kinematics map:
\begin{gather}
T = f (q) \text{ forward kinematics }\\
\mathcal{V} = J \dot{q}\text{ inverse kinematics }\\
\mathcal{A} = \dot{J} \dot{q} + J \ddot{q}
\end{gather}
This means $ \dot{J}  $ has a relation to acceleration, which we can use in computation
via Pinocchio.
\begin{gather}
\frac{d J (q (t))}{dt} =
\frac{\partial J}{\partial q} \dot{q} \\
= \left[ 
\frac{\partial J}{\partial q_{ 0 }} \dot{q} \dots \frac{\partial J}{\partial q_{ n }} \dot{q}
\right] 
\end{gather}
where $ H = \frac{\partial J}{\partial q}  \in \mathbb{R}^{d \times d \times n}$
is the forward kinematics Hessian.
We can't conveniently compute this directly.
However, Pinocchio offers the possibility to conveniently compute $ \dot{J}  $.
Looking at the last expression,
we can get $ H  $ by setting $ \dot{q} = e_{ i }  $ for every joint $ i  $.

Finally, we have all the components in a classic quadratic cost. This yields
the canonical quadratic cost $ x^{ T }Px + q^{ T }x  $ as follows:
\begin{gather}
		P = \mathcal{J}^{ T } \mathcal{J} + \lambda  I \\
		q = -2 * vec(L)^{  T} \mathcal{J}
\end{gather}
where we also threw in Tikhonov damping as standard practise 
($ \mathcal{J}  $ needs to be inverted within the QP solver, same story
as when you do a pseudoinverse.)

\section{Path-tracking redundancy resolution via path-induced manipulability tracking}
\subsection{How to know which manipulability ellipsoid to go to?}
In the paper, manipulability ellipsoids are learned from demonstration.
This is nifty because the demonstrator can have a different embodiment ---
the manipulability ellipsoid stays the same.

However, given a fully specified task such as door opening, one should
conceivably be able to construct an appropriate manipulability ellipsoid.
While it is obvious that, in the kinematic manipulability case,
the ellipsoid should be elongated along the direction of motion,
the quantitative question of how much remains unanswered.
The same goes for singularity avoiding strategies.

In this text we explore some options.


\printbibliography

\end{document}
